\section{Conclusion}

The 50 most costly freight bottlenecks in the US highway network impose an estimated \$22.7 billion in annual congestion costs. T2 connectors account for 64\% of bottleneck locations, but T1 Primary Arteries account for 62\% of total cost --- a cascade multiplier of 1.73 that reflects the national scope of T1 corridor disruptions versus the regional scope of T2 congestion.

The analysis confirms three findings from Paper A.1:

\textbf{The congestion-stress paradox holds}: aggregate-score T1 corridors (I-285, I-4, I-880) dominate bottleneck count but underperform on per-location cost. Centrality-adjusted T1 corridors (I-75, I-95, I-80) impose greater costs per bottleneck location through cascade effects.

\textbf{I-40 is exceptional}: the one T1 corridor where our analysis predicted no capacity emergency has zero ATRI top-50 locations. The convergence between ROUTE's model and ATRI's revealed-preference data provides external validation for the tier classification.

\textbf{Weather bottlenecks are a distinct category}: Donner Pass (\$1.6B/yr) is the 11th most expensive freight bottleneck in the US when episodic closures are included. The Donner freight tunnel is the highest-NPV single infrastructure investment in the I2.0 program.

The I2.0 investment sequencing that follows from this analysis:
\begin{enumerate}
\item Donner Pass freight tunnel (\$4B, 2.5yr payback)
\item T1 managed freight lanes on highest-bottleneck-cost corridors: I-95 Northeast, I-75 Atlanta, I-80 Bay Area
\item T2 relief in Atlanta (I-285 managed lanes / T1 upgrade) — simultaneously, to prevent I-285 congestion from undermining I-75/I-85 T1 improvements
\item Snoqualmie alternate alignment (\$3B, similar payback to Donner)
\end{enumerate}

The bottleneck data and the coverage gap data (Paper B.1) together define the full I2.0 investment geography: where the system is absent (coverage gaps) and where it is overwhelmed (bottlenecks). The missing link corridors (Northern Tier, I-69, I-3) address the former; the managed lane and diamond investments address the latter. Interstate 2.0 requires both.
